\section{Goal-generated ideal closure for nonlinear induction}
\label{sec:ideal}

\subsection{The multiplier certificate}

Keep the transition-system input of Section~\ref{sec:vector}, but now allow the entries of $H_a$ and $C$ to be polynomials in the state. The certificate conditions become
\begin{align}
 b_i(s_0)&=0 &&(1\le i\le r),\label{eq:ideal-init}\\
 b_i(F_a(X))&=\sum_{j=1}^r h_{aij}(X)b_j(X),\label{eq:ideal-step}\\
 g_\ell(X)&=\sum_{j=1}^r c_{\ell j}(X)b_j(X).
 \label{eq:ideal-target}
\end{align}
No multiplier has a sign condition. These are equality proofs, not positivity certificates. The acceptance language is deliberately just polynomial addition, multiplication, substitution, and rational equality.

\begin{theorem}[Ideal closure certificate]
Conditions \eqref{eq:ideal-init}--\eqref{eq:ideal-target} imply \eqref{eq:allwords}.
\end{theorem}
\begin{proof}
Assume every $b_j(s)$ is zero. Evaluating \eqref{eq:ideal-step} at $s$ gives $b_i(F_a(s))=0$, because every summand contains a zero factor. The initial condition and induction show that every basis polynomial vanishes at every reachable state. Evaluation of \eqref{eq:ideal-target} proves the target equalities.
\end{proof}

This proof does not use Gr\"obner bases, ideal membership algorithms, algebraic closure of the coefficient field, or a completeness theorem. Those enter search only. It works over any commutative ring in which the displayed coefficient and substitution identities are valid; the implemented rational coefficient representation imposes its own field-embedding obligations on a Lean frontend.

\subsection{An algorithm that avoids a bilinear synthesis problem}

Write $R=\Q[X_1,\ldots,X_d]$. The ideal generated by $b_1,\ldots,b_r$ is
\[
 \langle b_1,\ldots,b_r\rangle
 =\left\{\sum_j h_jb_j:h_j\in R\right\}.
\]
Instead of choosing an unknown template for every $b_j$ and every $h_j$, construct the $b_j$ from actual pullbacks of the goal. When a candidate $q$ is not in the current ideal, append $q$ itself as a new generator. Membership is decided during search by an untrusted exact Gr\"obner computation.

\begin{lstlisting}
B := []; J := zero ideal; work := original targets with origins
while work is nonempty:
    (q, goal, word) := pop_front(work)
    enforce resource limits
    if q(s0) != 0: return checked concrete word counterexample
    if q belongs to J: continue
    append the original polynomial q to B
    recompute a search-only Groebner basis for J = ideal(B)
    enqueue q composed with each transition, retaining word order
freeze B
for each required pullback and each target:
    reconstruct multipliers relative to the frozen B
    if reconstruction exceeds its separate budget: return UNKNOWN
return the explicit identities only after independent replay
\end{lstlisting}

The Gr\"obner basis is \emph{not} substituted for the retained generators. It is an index for membership queries. Retaining raw pullbacks makes a negative result interpretable as an execution of the original target and avoids making positive replay depend on a proof of the whole Gr\"obner algorithm.

\paragraph{Why fixed-generator reconstruction is linear.}
Once $B$ is frozen, choose a multiplier degree $e$ and write
\[
 h_j(X)=\sum_{|\alpha|\le e}u_{j\alpha}X^\alpha.
\]
The equation $q=\sum_jh_jb_j$ is linear in the unknown rational numbers $u_{j\alpha}$. Every coefficient of each $b_j$ is already fixed. The prototype first tries ordinary multivariate division by the raw generator list, then solves the exact coefficient system for increasing $e$ within the cap. An unsuccessful solve at one degree is not evidence that ideal membership is false.

A production producer could instead track transformation vectors while computing a Gr\"obner basis. If $G_k=\sum_j T_{kj}b_j$ and a reduction gives $q=\sum_k v_kG_k$, then $h_j=\sum_kv_kT_{kj}$. The same final identity checker suffices. The current package implements the fixed-generator linear reconstruction, not that faster provenance-carrying Gr\"obner backend.

\subsection{Why the unbounded algorithm terminates}

Define the pullback-generated ideal
\[
 J_* = \big\langle g_\ell\circ F_w:\ell\le t,\ w\in A^*\big\rangle.
\]
Every insertion strictly increases the current ideal. Unlike the vector space of all pullbacks, $R$ is Noetherian: every ascending chain of ideals stabilizes.

\begin{theorem}[Unbounded ideal discovery]
For finitely many total rational polynomial transitions, finitely many polynomial targets, and a fixed rational initial state, the uncapped discovery algorithm terminates. If multiplier reconstruction is allowed increasing finite degree bounds without a global cutoff, it returns either a certificate proving \eqref{eq:allwords} or a concrete execution refuting it.
\end{theorem}
\begin{proof}
Let $J_0\subsetneq J_1\subsetneq\cdots$ be the ideals after successful insertions. The Hilbert basis theorem forbids infinitely many strict inclusions. Each insertion creates only $|A|$ new work items, so after finitely many insertions the worklist also becomes empty.

Every retained generator is an actual target pullback, so its nonzero evaluation at the initial state yields a concrete counterexample. Suppose no such counterexample was returned. Every retained generator vanishes at $s_0$.

At exhaustion the ideal $J$ contains every target and every pullback of each retained generator. Since substitution is a ring homomorphism,
\[
 F_a^*\left(\sum_jh_jb_j\right)
 =\sum_j(F_a^*h_j)(F_a^*b_j)\in J.
\]
Thus $J$ is closed under every pullback and contains all of $J_*$. Conversely every inserted generator belongs to $J_*$, so $J=J_*$. The required step and target identities therefore have polynomial multipliers. Each such representation has some finite degree; increasing-degree coefficient reconstruction eventually finds one. The initial, step, and target conditions then give a positive certificate.

Finally, if a counterexample $w$ existed after positive closure, $g_\ell\circ F_w\in J_*$ would evaluate to zero at $s_0$, since every generator does. This contradicts its definition as a counterexample.
\end{proof}

\begin{remark}[Precisely what was decided]
This is a decision result for equality observations on the closed, unguarded polynomial system just defined. It is not a decision result for arbitrary program verification, real inequalities, unrestricted guards, termination, general Lean goals, or the existence of a term inhabiting a dependent type. The executed implementation is resource-bounded and therefore does not implement an unconditional decision interface.
\end{remark}

The Noetherian proof is a termination argument, not a practical degree or runtime bound. A single Gr\"obner calculation may dominate runtime and memory. The current synchronous SymPy calls do not provide fine-grained cancellation. Production deployment must run the producer under an external process deadline and memory limit; a timeout is \unknown{}. This caveat is not removed by the existence of a mathematical termination theorem.

\subsection{Why a radical ideal is unnecessary}

It is tempting to replace $J$ by its radical because only zero sets matter. That can shrink the generating set, but it creates new proof obligations. A radical-membership claim $q^m\in J$ only yields $q=0$ in a setting with appropriate reducedness or integral-domain properties. The direct multiplier certificate makes no such inference.

For this task ordinary ideals already yield the termination theorem. The package deliberately avoids radical computation. A later producer may use radical information to propose simpler generators, but it must add a replay rule that proves the relevant root-of-zero implication in the source coefficient domain. This is a concrete place where a search optimization would otherwise silently strengthen the logic of the checker.

\subsection{Worked nonlinear separation: repeated squaring}

Let $F(x,y)=(x^2,y^2)$, $s_0=(2,2)$, and $g=x-y$. The vector worker reaches its degree ceiling after seeing degree $16$, having retained four independent orbit polynomials. The ideal worker retains only $b_1=x-y$ and supplies
\[
 b_1(F(x,y))=(x+y)b_1(x,y),\qquad g=b_1.
\]
The independent checker expands both sides and checks $b_1(2,2)=0$. The fact that values become enormous under iteration is irrelevant to positive checking: no long trajectory is evaluated.

The generated nonlinear corpus uses two synchronized scalar maps, including powers two through four, linear terms, and constants. It is intentionally a small structural family with known equivalence. Its value is to test multiplier generation across different symbolic updates; it is not an unbiased measure of all nonlinear verification problems.

\subsection{Worked generator discovery beyond one equation}

Take $F(x,y)=(y,x+y)$, $s_0=(0,0)$, and the target $b_0=xy$. The first pullback is
\[
 b_1=xy+y^2,
\]
which does not belong to $\langle xy\rangle$. Its pullback is
\[
 b_2=x^2+3xy+2y^2,
\]
which adds another generator. These three span the quadratic monomials as a vector space and therefore generate the ideal $(x,y)^2$. The next pullback satisfies
\[
 b_2\circ F=-b_0+2b_1+2b_2.
\]
Consequently the generated joint zero condition is inductive. This illustrates the distinction between simply verifying a user-supplied ideal invariant and discovering the stronger invariant needed by the goal.

The initial state in this example remains fixed, making the theorem itself easy. That does not invalidate the generator-growth test, but it does prevent the case from supporting a claim of superiority to existing Lean automation. The package keeps that distinction explicit.

\subsection{Guards, control flow, and initial assumptions}

For guarded transitions, an unconditional multiplier identity is still a sufficient certificate: it proves preservation even when the guard is discarded. The proof may become harder or impossible after that abstraction. A counterexample word in the abstraction must not refute the source statement until every guard is checked on the resulting concrete states.

A practical extension for finite control-flow graphs uses one generator family $B_v$ per program point. An edge $u\xrightarrow{F}v$ requires
\[
 B_v(F(X))=H_{uv}(X)B_u(X).
\]
Targets seed their program points, and work propagates backward across incoming edges. A finite tuple of ideal lattices over Noetherian rings still admits the same ascending-chain argument. Initial evaluation is needed only at entry points. This extension is specified here, not implemented in the supplied single-location worker.

For an initial state parameterized polynomially by $z$, replace the scalar initial checks with identities $b_i(s_0(z))=0$. That proves the result for every rational parameter value. For an initial \emph{assumption} ideal $K=\langle k_1,\ldots,k_m\rangle$, a sufficient initial certificate is $b_i=\sum_jv_{ij}k_j$. This need not be complete for every semantic initial zero set unless further algebraic reasoning is introduced; the implementation must not claim completeness merely from a sufficient ideal-membership test.

\subsection{A further extension: arbitrary polynomial input values}
\label{sec:inputcoeff}

Finite alphabets are useful, but many folds consume integers or rationals. Suppose a transition is polynomial in both state $X$ and an arbitrary input $Z$. Expand
\[
 b_i(F(X,Z))=\sum_\alpha q_{i\alpha}(X)Z^\alpha.
\]
A sufficient closure certificate expresses every coefficient $q_{i\alpha}$ in the state ideal. Equivalently, supply multipliers $h_{ij}(X,Z)$ with
\[
 b_i(F(X,Z))=\sum_jh_{ij}(X,Z)b_j(X).
\]
This establishes preservation for \emph{all} input values, and then induction establishes a contract for lists of arbitrary length and arbitrary rational entries.

For discovery, enqueue the finitely many coefficient polynomials of each pullback, rather than treating infinitely many input values as separate transition symbols. Each insertion still strictly increases a state ideal, so Noetherianity gives a positive closure route. However, the raw-word counterexample mechanism must change: a coefficient polynomial is not itself the observation along a concrete input word. One must retain the parameterized word and reconstruct actual rational inputs, or restrict this extension to positive evidence. The package does not implement this extension and does not reuse the finite-alphabet negative certificate unmodified.

This coefficientwise extension is particularly relevant to Forge's relationship with Leant. A universal assertion about a polynomial accumulator can be discharged through a finite symbolic induction invariant even when the input element type is infinite. Type-directed synthesis can propose the program; this worker can prove a scoped behavioral contract, without assuming that type correctness implies the intended behavior.
